\section{Implementation}
\label{sec:impl}

\td{Built on the DaCe dataflow compiler + its Fortran frontend (\texttt{f2dace}):
\texttt{flang -fc1 -emit-hlfir} $\to$ MLIR normalisation $\to$ SDFG $\to$ specialise
$\to$ regenerate a \texttt{bind(c)} wrapper. List the passes: glacial recovery,
\texttt{propagate\_if\_cond}, \texttt{LiftTrivialIf}, \texttt{compress\_indices}
(\#3), \texttt{\_monomorphize\_solver} (\#4), the covering-set linker.}
\td{INTEGRATION-AS-A-PASS (reviewer-facing): the value$\to$type devirtualization
belongs at the FIR/HLFIR level, where \texttt{fir.dispatch} and the type descriptor
are explicit and the dynamic type is recoverable --- by LLVM IR the type info has
been erased into a descriptor-table indirect call. So the approach is a
\emph{FIR-level pass}, prototyped via our DaCe HLFIR consumption and a natural Flang
upstream; it is not DaCe-specific. The scalar bridge (C1) is an IPSCCP-seeding pass
(seed glacial globals with recovered constants). Frame whole-program
devirtualization for C++ as the precedent we extend to a new dispatch model and a
new driver (config value, not the type hierarchy).}
\td{Multi-variant linking: per-variant \texttt{ld -r} + \texttt{objcopy
--localize-symbols} on colliders; public per-SDFG entry points stay global for
name-based startup dispatch.}
\td{Serde determinism (DaCe PR \#2392) underpins reliable SDFG round-tripping.}
